\begin{theorem}[Lee--Yang 根谱的逐根可加性：纤维条件 KL 的 Bernoulli 分解]\label{thm:pom-fiber-leyang-kl-bernoulli-decomposition}
固定 \(m\ge 1\) 与 \(x\in X_m\)，并沿用定理 \ref{thm:pom-fiber-biased-posterior-poisson-binomial} 的 Lee--Yang 根分解
\[
Z_x(z)=C_x\prod_{j=1}^{d_x}\left(1+\frac{z}{\alpha_{x,j}}\right),
\qquad
\alpha_{x,j}>0,
\qquad
q_{x,j}(t):=\frac{t}{t+\alpha_{x,j}}.
\]
则对任意 \(p,q\in(0,1)\) 与 \(t_p=p/(1-p)\)、\(t_q=q/(1-q)\) 有严格恒等式
\[
D_{\mathrm{KL}}\!\left(\mu_{m,p}(\cdot\mid x)\,\middle\|\,\mu_{m,q}(\cdot\mid x)\right)
=\sum_{j=1}^{d_x}
D_{\mathrm{KL}}\!\left(\mathrm{Bern}(q_{x,j}(t_p))\,\middle\|\,\mathrm{Bern}(q_{x,j}(t_q))\right).
\]
\end{theorem}
\begin{proof}
对任意 \(\alpha>0\) 定义 \(q(t)=t/(t+\alpha)\)。
直接计算 Bernoulli 相对熵得
\[
D_{\mathrm{KL}}\!\left(\mathrm{Bern}(q(t_p))\,\middle\|\,\mathrm{Bern}(q(t_q))\right)
=q(t_p)\log\frac{t_p}{t_q}-\log\frac{\alpha+t_p}{\alpha+t_q}.
\]
对 \(\alpha=\alpha_{x,j}\) 求和并注意
\[
\sum_{j=1}^{d_x}q_{x,j}(t_p)=\EE_{\mu_{m,p}(\cdot\mid x)}[R_{x,p}]
\]
（定理 \ref{thm:pom-fiber-biased-posterior-poisson-binomial}）以及
\[
\log\frac{Z_x(t_p)}{Z_x(t_q)}
=\sum_{j=1}^{d_x}\log\frac{\alpha_{x,j}+t_p}{\alpha_{x,j}+t_q}
\]
（由根分解取对数）即得
\[
\sum_{j=1}^{d_x}
D_{\mathrm{KL}}\!\left(\mathrm{Bern}(q_{x,j}(t_p))\,\middle\|\,\mathrm{Bern}(q_{x,j}(t_q))\right)
=\EE[R_{x,p}\mid x]\log\frac{t_p}{t_q}-\log\frac{Z_x(t_p)}{Z_x(t_q)}.
\]
与定理 \ref{thm:pom-fiber-biased-conditional-kl-chain} 的闭式 \eqref{eq:pom-fiber-biased-conditional-kl-closed} 一致，故结论成立。证毕。
\end{proof}

\begin{theorem}[Lee--Yang 根谱的逐根可加性：纤维条件 R\'enyi 散度分解]\label{thm:pom-fiber-leyang-renyi-bernoulli-decomposition}
固定 \(m\ge 1\) 与 \(x\in X_m\)，并沿用定理 \ref{thm:pom-fiber-biased-posterior-poisson-binomial} 的记号。
对 \(s>0\)、\(s\neq 1\)，定义有限样本空间上的 R\'enyi 散度
\[
D_s(P\|Q):=\frac{1}{s-1}\log\sum_{\omega}P(\omega)^sQ(\omega)^{1-s}.
\]
则对任意 \(p,q\in(0,1)\) 有恒等式
\[
D_s\!\left(\mu_{m,p}(\cdot\mid x)\,\middle\|\,\mu_{m,q}(\cdot\mid x)\right)
=\sum_{j=1}^{d_x}
D_s\!\left(\mathrm{Bern}(q_{x,j}(t_p))\,\middle\|\,\mathrm{Bern}(q_{x,j}(t_q))\right).
\]
并且当 \(s\to 1\) 时该分解连续退化为定理 \ref{thm:pom-fiber-leyang-kl-bernoulli-decomposition} 的 KL 分解。
\end{theorem}
\begin{proof}
由 \(\mu_{m,p}(\omega\mid x)=t_p^{r(\omega)}/Z_x(t_p)\) 得
\[
\sum_{\omega\in\Fold_m^{-1}(x)}\mu_{m,p}(\omega\mid x)^s\mu_{m,q}(\omega\mid x)^{1-s}
=\frac{Z_x(t_p^s t_q^{1-s})}{Z_x(t_p)^s Z_x(t_q)^{1-s}}.
\]
将根分解代入并逐因子化简，得到
\[
\frac{Z_x(t_p^s t_q^{1-s})}{Z_x(t_p)^s Z_x(t_q)^{1-s}}
=\prod_{j=1}^{d_x}
\frac{\alpha_{x,j}+t_p^s t_q^{1-s}}{(\alpha_{x,j}+t_p)^s(\alpha_{x,j}+t_q)^{1-s}}.
\]
另一方面，对 \(q(t)=t/(t+\alpha)\) 有恒等式
\[
q(t_p)^s q(t_q)^{1-s}+\bigl(1-q(t_p)\bigr)^s\bigl(1-q(t_q)\bigr)^{1-s}
=\frac{\alpha+t_p^s t_q^{1-s}}{(\alpha+t_p)^s(\alpha+t_q)^{1-s}}.
\]
因此取对数并除以 \(s-1\) 得到逐根相加的 R\'enyi 分解。
极限 \(s\to 1\) 的连续性在有限样本空间上是标准事实，故可退化得到 KL 情形。证毕。
\end{proof}

\begin{corollary}[条件 Chernoff 信息的严格凸谱表达与唯一极小点]\label{cor:pom-fiber-leyang-chernoff-unique-minimizer}
固定 \(m\ge 1\) 与 \(x\in X_m\)，取 \(p,q\in(0,1)\)。
定义条件 Chernoff 信息
\[
\mathcal{C}_x(p,q)
:=
-\log\!\left(\inf_{s\in[0,1]}\sum_{\omega}\mu_{m,p}(\omega\mid x)^s\,\mu_{m,q}(\omega\mid x)^{1-s}\right).
\]
则有变分表示
\[
\mathcal{C}_x(p,q)
=\max_{s\in[0,1]}
\Bigl(
s\log Z_x(t_p)+(1-s)\log Z_x(t_q)-\log Z_x(t_p^s t_q^{1-s})
\Bigr),
\]
并且函数
\[
s\ \longmapsto\ \log Z_x(t_p^s t_q^{1-s})-s\log Z_x(t_p)-(1-s)\log Z_x(t_q)
\]
在 \((0,1)\) 上凸；
若 \(p\neq q\) 且 \(Z_x\) 非常值，则其在 \((0,1)\) 上严格凸，从而 \(\mathcal{C}_x(p,q)\) 的极小点唯一（等价地，上述对偶极大点唯一）。
进一步，借助 Lee--Yang 根分解可写成
\[
\mathcal{C}_x(p,q)
=\max_{s\in[0,1]}\sum_{j=1}^{d_x}
\Bigl(
s\log(\alpha_{x,j}+t_p)+(1-s)\log(\alpha_{x,j}+t_q)-\log(\alpha_{x,j}+t_p^s t_q^{1-s})
\Bigr).
\]
\end{corollary}
\begin{proof}
由 R\'enyi “幂平均”恒等式（定理 \ref{thm:pom-fiber-leyang-renyi-bernoulli-decomposition} 的证明）得
\(
\sum_{\omega}\mu_{m,p}(\omega\mid x)^s\mu_{m,q}(\omega\mid x)^{1-s}
=Z_x(t_p^s t_q^{1-s})/(Z_x(t_p)^s Z_x(t_q)^{1-s})
\)。
代入定义并取 \(-\log\)，得到所示对偶极大形式。
\par
令 \(A(\eta):=\log Z_x(e^\eta)\) 且 \(\eta_p=\log t_p\)、\(\eta_q=\log t_q\)，并记 \(\eta(s)=s\eta_p+(1-s)\eta_q\)。
则上式中的括号项恰为
\(
A(\eta(s))-sA(\eta_p)-(1-s)A(\eta_q)
\)。
由指数族性质 \(A''(\eta)=\mathrm{Var}(R_{x,p}\mid x)\ge 0\) 可得该 Jensen 差的二阶导数为
\(
\frac{\dd^2}{\dd s^2}\bigl(A(\eta(s))-sA(\eta_p)-(1-s)A(\eta_q)\bigr)
=(\eta_p-\eta_q)^2A''(\eta(s))
\)，
从而该函数在 \((0,1)\) 上凸；并且当 \(p\neq q\) 且 \(Z_x\) 非常值时，\(A''(\eta(s))>0\) 导出严格凸性并从而极小点唯一。
最后一式由根分解把 \(\log Z_x\) 写成求和后逐项展开即得。证毕。
\end{proof}

\begin{proposition}[由条件 KL 曲线反演 Lee--Yang 根谱]\label{prop:pom-fiber-leyang-roots-identifiable-from-kl-curve}
固定 \(m\ge 1\) 与 \(x\in X_m\)，并令
\[
D_x(t\,\|\,t_0):=D_{\mathrm{KL}}\!\left(\mu_{m,p(t)}(\cdot\mid x)\,\middle\|\,\mu_{m,p(t_0)}(\cdot\mid x)\right),
\qquad
p(t):=\frac{t}{1+t}.
\]
若函数 \(t\mapsto D_x(t\,\|\,t_0)\) 在某个开区间 \(I\subset(0,\infty)\) 上可观测，则 \(\log Z_x(t)\) 在 \(I\) 上可由 \(D_x(\cdot\,\|\,t_0)\) 唯一恢复至多差一个与 \(t\) 无关的加性常数。
从而 \(Z_x\) 的 Lee--Yang 根多重集合 \(\{\alpha_{x,j}\}_{j=1}^{d_x}\) 被 \(D_x(\cdot\,\|\,t_0)\) 唯一确定。
\end{proposition}
\begin{proof}
令自然参数 \(\eta=\log t\)，并定义对数配分函数 \(A(\eta):=\log Z_x(e^\eta)\)。
由定理 \ref{thm:pom-fiber-biased-conditional-kl-chain} 的闭式 \eqref{eq:pom-fiber-biased-conditional-kl-closed} 可整理为
\[
D_x(e^\eta\,\|\,e^{\eta_0})
=A'(\eta)(\eta-\eta_0)-\bigl(A(\eta)-A(\eta_0)\bigr),
\qquad \eta_0:=\log t_0.
\]
对 \(\eta\) 求导得
\[
\frac{\dd}{\dd\eta}D_x(e^\eta\,\|\,e^{\eta_0})
=A''(\eta)(\eta-\eta_0),
\]
从而在 \(\eta\neq\eta_0\) 处有
\(
A''(\eta)=\bigl(\frac{\dd}{\dd\eta}D_x(e^\eta\,\|\,e^{\eta_0})\bigr)/(\eta-\eta_0)
\)。
因此 \(D_x(\cdot\,\|\,t_0)\) 在开区间上确定 \(A''\)，进而两次积分确定 \(A\) 在该区间上至多差一个仿射函数。
但 \(A(\eta)=\log Z_x(e^\eta)\) 在 \(\eta\to-\infty\) 时有界（因 \(Z_x(0)=C_x\in(0,\infty)\)），迫使该仿射不确定性无斜率项，仅余加性常数。
最后，由 \(Z_x\) 在开区间上的值（至多差乘性常数）即可唯一决定其根集合；乘性常数不影响根谱，故 \(\{\alpha_{x,j}\}\) 被唯一确定。证毕。
\end{proof}

